\section{OPD 使用方式与误差传递}
\label{sec:kl-opd-roles}
\subsection{Score baseline 是积分为零的控制}

本小节默认 $\rho=p$。vOPD 在总定理中对应 $f_v=r_vs_v$、$a_v=c(x)s_v$。由于 $I_p(a)=c(x)I_p(s)=0$，它直接成为 $(r_A-c)s_A$。对其他覆盖残差支持的采样分布，正确形式是 $(p(A)/\rho(A))(r_A-c)s_A$。因此 baseline 是否对 KL 数值无偏，不是这个构造的无偏条件；控制积分为零才是条件。下面求出的最优 $c$ 就是 Gram 定理的单列控制解。


Sampled-token estimator 为
\[
 \widehat g^{\mathrm{sample}}=r_As_A.
\]
它无偏，且
\begin{equation}
 \Var(\widehat g^{\mathrm{sample}})
 =\E_p[r_A^2\lVert s_A\rVert_2^2]-\lVert g\rVert_2^2.
 \label{eq:chapter3-sample-grad-var}
\end{equation}

对任意在抽样前固定的 scalar $c(x)$，vOPD 型 estimator 为
\begin{equation}
 \widehat g^{\mathrm{vOPD}}_c=(r_A-c)s_A.
 \label{eq:chapter3-vopd}
\end{equation}

\begin{proposition}[固定 scalar baseline 的无偏性与最优值]
\label{prop:chapter3-vopd}
对任意 action-independent $c(x)$，$\widehat g^{\mathrm{vOPD}}_c$ 无偏。其二阶矩为
\[
 \E\lVert\widehat g^{\mathrm{vOPD}}_c\rVert_2^2
 =\E_p[(r_A-c)^2\lVert s_A\rVert_2^2],
\]
当 $\E_p\lVert s_A\rVert^2>0$ 时，使该二阶矩最小的 scalar 是
\begin{equation}
 c^*=\frac{\E_p[r_A\lVert s_A\rVert_2^2]}{\E_p[\lVert s_A\rVert_2^2]}.
 \label{eq:chapter3-optimal-score-baseline}
\end{equation}
一般 $c^*\neq K$，因此用 KL 数值 $K$ 作为 vOPD baseline 是合法的中心化，但不必是 score-gradient 方差最优的 baseline。
若 $\E_p\lVert s_A\rVert^2=0$，score 在 $p$ 下几乎处处为零，所有 baseline 对应的梯度均为零。
\end{proposition}

\begin{proof}
\[
 \E_p[(r_A-c)s_A]=\E_p[r_As_A]-c\E_p[s_A]=g.
\]
对二阶矩关于 $c$ 求导并令其为零，得到式~\eqref{eq:chapter3-optimal-score-baseline}。
\end{proof}

进一步记 $D_s=\E_p\lVert s_A\rVert^2>0$，展开平方并配方，得到精确的方差表达式
\begin{align*}
 \Var(\widehat g_c^{\mathrm{vOPD}})
 &=\E_p[r_A^2\lVert s_A\rVert^2]
   -2c\E_p[r_A\lVert s_A\rVert^2]+c^2D_s-\lVert g\rVert^2\\
 &=\Var(\widehat g_{c^*}^{\mathrm{vOPD}})+D_s(c-c^*)^2.
\end{align*}
这给出任意固定 KL proxy 作为 baseline 的代价，而无需假定该 proxy 对 $K$ 无偏。

\begin{proposition}[独立随机 KL baseline 的误差传递]
令 $C$ 与当前动作 $A$ 条件独立，$m=\E[C]$、$\sigma_C^2=\Var(C)$，且反传时停止 $C$ 的梯度。则 $(r_A-C)s_A$ 对 $g$ 无偏，且
\[
 \Var((r_A-C)s_A)
 =\Var((r_A-m)s_A)+D_s\sigma_C^2.
\]
若写 $m=K+\delta$，则
\[
 \Var((r_A-C)s_A)
 =\Var(\widehat g_{c^*}^{\mathrm{vOPD}})
  +D_s\{(K+\delta-c^*)^2+\sigma_C^2\}.
\]
\end{proposition}

\begin{proof}
条件于 $A$，$\E_C[(r_A-C)^2]=(r_A-m)^2+\sigma_C^2$。乘以 $\lVert s_A\rVert^2$ 后取 $A$ 的期望，并减去相同的 $\lVert g\rVert^2$，即可得到结论。
\end{proof}

所以在 vOPD 中只改进后面的 KL 数值估计时，应比较 scalar baseline 的均值与方差；当两个 baseline 均对 $K$ 无偏且独立时，梯度方差之差恰好是 $D_s$ 乘以 scalar 方差之差。ResOPD 的向量控制方差结论不能直接用于这个比较。

如果 baseline 是与当前 action 共同随机化的 $C_A$，则上述证明失效：偏差为
\begin{equation}
 \operatorname{Bias}(\widehat g_C)=-\E[C_As_A].
 \label{eq:chapter3-random-baseline-bias}
\end{equation}
因此随机 baseline 保持无偏的充要条件是 $\E[C_As_A]=0$。独立样本或 leave-one-out/cross-fitting 是保证这个条件的构造方式；仅有 $\E[C_A]=K$ 不够。

同样本 K1 baseline 提供一个极端反例：取 $C_A=r_A$，则 $(r_A-C_A)s_A=0$，偏差恰好为 $-g$，尽管 $\E[C_A]=K$。数值无偏不能消除 baseline 与 score 的相关性。

\subsection{旧 rollout 下的最优 baseline}

若 $A\sim\rho$，且 $\rho$ 覆盖 $p$ 的支持，令 $w_A=p(A)/\rho(A)$。控制 $c s_A$ 的正确修正仍然无偏，因为
\[
 \E_\rho[w_A(r_A-c)s_A]=I_p(rs)-cI_p(s)=g.
\]
但是其二阶矩使用 $p^2/\rho$ 权重。定义
\[
 D_{s,\rho}=\sum_v\frac{p(v)^2}{\rho(v)}\|s_v\|^2.
\]
若该值为正，对二阶矩关于 $c$ 配方得到
\[
 c_\rho^*=
 \frac{\sum_v p(v)^2\rho(v)^{-1}r_v\|s_v\|^2}{D_{s,\rho}},
\]
\[
 \Var_\rho(w_A(r_A-c)s_A)
 =\Var_\rho(w_A(r_A-c_\rho^*)s_A)
  +D_{s,\rho}(c-c_\rho^*)^2.
\]
这直接来自 Gram 定理取控制列 $w_As_A$。若 $D_{s,\rho}=0$，所有 score 均为零，baseline 不影响结果。

对与当前动作独立的随机 baseline $C$，令 $m=\E C$、$\sigma_C^2=\Var C$。条件于动作展开平方，得到
\[
 \Var_\rho(w_A(r_A-C)s_A)
 =\Var_\rho(w_A(r_A-m)s_A)+D_{s,\rho}\sigma_C^2.
\]
因此两个独立、同均值的 KL baseline 的比较，只需比较其标量方差，梯度方差差为 $D_{s,\rho}$ 乘以该差；当前 on-policy 情形是 $\rho=p$ 的特例。若 baseline 与动作相关，无偏性的条件也相应变成 $\E_\rho[w_AC_As_A]=0$。

\subsection{Sampled OPD 与 vOPD}

固定 teacher 作为 $q$，sampled OPD 取停止系数 $r_A$，其最小化梯度为 $r_As_A$；若使用 reward 最大化记号，reward 是 $-r_A$，两种表述符号相反而更新方向相同。vOPD 取 $(r_A-c(x))s_A$，其中 $c(x)$ 可取完整 $K$、Tail-aware 下界或独立随机 KL baseline。第三章已经证明：action-independent baseline 的数值偏差不产生梯度偏差；同 action 的无偏 KL 估计反而可以产生梯度偏差。

本仓库生产 OPD baseline 对 $q=\pi_T$ 使用 $H+\tau_p\log(\tau_p/\tau_q)$，并在梯度系数中停止。它是完整 KL 的有偏数值下界，但作为 action-independent baseline，在正确采样测度或正确重要性校正下保持 OPD 梯度无偏。若 rollout 来自冻结 $\pi_k$，则此处 $\rho=\pi_k$，需要 $p/\rho$；baseline 本身使用 teacher 作为 $q$，二者不可混用。

\subsection{Mirror 目标中的 student--anchor KL}

对 MDPO，令本节参考 $q=\pi_k$，固定 reward 为前两章的 $r_k(x,v)$。负 mirror 目标为
\[
 \mathcal L_{\mathrm{MDPO}}
 =\eta^{-1}K-\E_p[r_k(x,A)].
\]
逐项求导并利用 score 均值为零，得到
\[
 \nabla\mathcal L_{\mathrm{MDPO}}
 =\E_p[u_A^{\mathrm{MDPO}}s_A],\qquad
 u_v^{\mathrm{MDPO}}=\eta^{-1}r_v-r_k(x,v).
\]
在冻结 rollout $\rho=\pi_k$ 下，无偏估计为
\[
 \widehat g_{\mathrm{MDPO}}=\frac{p(A)}{\rho(A)}u_A^{\mathrm{MDPO}}s_A,\qquad
 \Var_\rho(\widehat g_{\mathrm{MDPO}})
 =\sum_v\frac{p(v)^2}{\rho(v)}(u_v^{\mathrm{MDPO}})^2\lVert s_v\rVert^2
  -\lVert\nabla\mathcal L_{\mathrm{MDPO}}\rVert^2.
\]
如果 $\lVert u_v^{\mathrm{MDPO}}s_v\rVert\leq C$，则二阶矩满足
\begin{equation}
 \E_\rho\lVert\widehat g_{\mathrm{MDPO}}\rVert^2
 \leq C^2\sum_v\frac{p(v)^2}{\rho(v)}
 =C^2\{1+\chi^2(p\Vert\rho)\}.
 \label{eq:mdpo-variance}
\end{equation}
这说明 stale rollout 与当前 student 分离时的重要性权重代价。

BPO 的 student center 是相反方向的 $\KL(\pi_k\Vert p)$，在本节 $q=\rho=\pi_k$ 下就是 $K^{\mathrm F}$；teacher center 是冻结的 $\KL(\pi_k\Vert\pi_T)$。PMD-Mean 的目标使用冻结 reward 中心而不含同样的 student KL center。它们的原目标就在 anchor 测度下定义，不能因 MDPO 需要 $p/\rho$，便给 Mean 或 BPO 机械乘上该比率。center 放入平方损失后还会出现同样本相关项，须使用第二章的有限样本与留一修正；一个 KL scalar estimator 的无偏性不替代整个目标梯度的证明。

\subsection{确定性代理的平方误差}

记 $M=\tau_p\ell_{\tail}$、$B=\tau_p\KL(p_{\tail}\Vert q_{\tail})$。Tail-aware 与截断的平方误差差为
\[
 (K-H)^2-\{K-(H+M)\}^2=(M+B)^2-B^2=M(M+2B).
\]
因此 $M\geq0$ 时 Tail-aware 的标量平方误差不大于截断；$M<0$ 时删除的质量项与形状项可抵消，零方差确定性 proxy 的排序不能只按是否保留更多项判断。

支持细分的 KL 单调性已经由第~\ref{sec:kl-partition} 节的一般分区定理证明；该性质控制确定性粗化偏差，不能推出 HETS 或 ResOPD 的抽样方差单调下降。

\subsection{Normalized 与 Truncated 的 logit 行为}

令 $p=\operatorname{softmax}(\theta)$，其 score 为 $s_{v,j}=\mathbf1\{v=j\}-p(j)$。Normalized Top-$J$ 只依赖头部相对 logits，故对 $j\notin S$ 有
\[
 \frac{\partial K^{\mathrm{norm}}_J}{\partial\theta_j}=0.
\]
截断数值的真实梯度则满足
\[
 \frac{\partial H}{\partial\theta_j}=-p(j)(H+P),\qquad j\notin S.
\]
这由 $\nabla H=\sum_{v\in S}p(v)(r_v+1)s_v$ 逐坐标代入得到。当 $H+P>0$ 时，梯度下降增加尾部 logits。停止系数的头部梯度 $h_g$ 省去了 $P$ 对应的边界项，但仍缺失完整尾部。因此这两种截断训练路径也必须分别报告。

\subsection{固定 prefix 误差与访问测度下误差}

令 $x\sim d$，条件偏差为 $\beta(x)=\E[\widehat K\mid x]-K(x)$，条件方差为 $V(x)$。则
\[
 \E_{x,A}[\widehat K-K]=\E_d[\beta(x)],\qquad
 \E_{x,A}[(\widehat K-K)^2]=\E_d[\beta(x)^2+V(x)].
\]
而原始 estimator 的总体方差是
\[
 \Var_{x,A}(\widehat K)
 =\E_d[V(x)]+\Var_d(K(x)+\beta(x)).
\]
这三个统计量分别回答总体有符号误差、逐 prefix 均方误差和 estimator 自身的总体波动。确定性 Tail-aware 有 $V(x)=0$，但 $\operatorname{RMSE}=\sqrt{\E_d[B(x)^2]}$，一般不等于 $|\E_d[-B(x)]|$。向量估计量将平方替换为范数平方、方差替换为协方差迹，结论保持相同形式。访问测度 $d$ 在不同 outer 轮次的变化由第四章单独处理。
